% ============================================================================
%  CNC302 — Лабораторийн ажлын гарын авлага
%  Тохиргоо (preamble). XeLaTeX эсвэл LuaLaTeX-ээр хөрвүүлнэ.
% ============================================================================

% ── Хуудасны хэмжээ ─────────────────────────────────────────────────────────
\usepackage[
  a4paper,
  top=22mm, bottom=22mm, left=24mm, right=20mm,
  headheight=18pt, footskip=12mm
]{geometry}

% ── Фонт (кирилл үсэг) ──────────────────────────────────────────────────────
\usepackage{fontspec}
\usepackage{amssymb}   % $\square$ (pandoc-ийн checkbox жагсаалт)
\defaultfontfeatures{Ligatures=TeX}
\setmainfont{DejaVu Serif}[
  Scale=0.94,
  BoldFont={DejaVu Serif Bold},
  ItalicFont={DejaVu Serif Italic},
  BoldItalicFont={DejaVu Serif Bold Italic}]
\setsansfont{DejaVu Sans}[Scale=0.90]
\setmonofont{DejaVu Sans Mono}[Scale=0.82]

% Монгол хэлэнд тусгай мөр таслалт байхгүй тул зөөлөн зохицуулалт хийнэ
\sloppy
\emergencystretch=3em
\hyphenpenalty=1000
\tolerance=2000

% ── Тусгай тэмдэгт (DejaVu Serif-д байхгүй, DejaVu Sans-аас авна) ──────────
\usepackage{newunicodechar}
\newfontfamily{\symfont}{DejaVu Sans}[Scale=0.90]
\newcommand{\sym}[1]{{\symfont #1}}
\newunicodechar{✓}{\sym{✓}}
\newunicodechar{✔}{\sym{✔}}
\newunicodechar{✗}{\sym{✗}}
\newunicodechar{✘}{\sym{✘}}
\newunicodechar{⚠}{\sym{⚠}}
\newunicodechar{⚙}{\sym{⚙}}
\newunicodechar{☰}{\sym{☰}}
\newunicodechar{◆}{\sym{◆}}
\newunicodechar{●}{\sym{●}}
\newunicodechar{▼}{\sym{▼}}
\newunicodechar{▲}{\sym{▲}}
\newunicodechar{✖}{\sym{✖}}
\newunicodechar{★}{\sym{★}}
\newunicodechar{↔}{\sym{↔}}
\newunicodechar{⌈}{\sym{⌈}}
\newunicodechar{⌉}{\sym{⌉}}
\newunicodechar{σ}{\sym{σ}}

% ── Өнгө ────────────────────────────────────────────────────────────────────
\usepackage{xcolor}
\definecolor{cncblue}{HTML}{1F3864}
\definecolor{cncmid}{HTML}{2E5496}
\definecolor{cnclight}{HTML}{D9E2F3}
\definecolor{cncgrey}{HTML}{F2F2F2}
\definecolor{cncrule}{HTML}{8EA9DB}
\definecolor{cncwarn}{HTML}{B45309}
\definecolor{cncok}{HTML}{15803D}
\definecolor{cncdanger}{HTML}{B91C1C}

% ── Хүснэгт ─────────────────────────────────────────────────────────────────
\usepackage{longtable}
\usepackage{booktabs}
\usepackage{array}
\usepackage{calc}
\usepackage{tabularx}
\usepackage{colortbl}
\setlength{\tabcolsep}{4pt}
\renewcommand{\arraystretch}{1.18}
\newcolumntype{L}[1]{>{\raggedright\arraybackslash}p{#1}}

% Урт хүснэгтийн толгойг давтахад "үргэлжлэл" гэж тэмдэглэнэ
\usepackage{caption}
\captionsetup{font=small,labelfont=bf,skip=4pt}

% ── График, диаграм ─────────────────────────────────────────────────────────
\usepackage{graphicx}
\usepackage{float}
% Зураг: хуудасны өргөнөөр, өндрийг хязгаарлаж, бичвэрийн дунд байрлуулна.
% Гарын авлага тул зургийг ДАВХАР хөвүүлэхгүй — яг тайлбарласан газраа ([H]).
\setkeys{Gin}{width=\linewidth,height=0.40\textheight,keepaspectratio}
\floatplacement{figure}{H}
% Тайлбар нь өөрөө "Зураг 1.1 — …" гэж эхэлдэг тул LaTeX-ийн дугаарыг нэмэхгүй
\captionsetup[figure]{labelformat=empty,font=small,skip=6pt}

% ── Жагсаалт ────────────────────────────────────────────────────────────────
\usepackage{enumitem}
\setlist{topsep=3pt,itemsep=1pt,parsep=1pt,leftmargin=*}
\setlist[itemize]{label=\textcolor{cncmid}{\textbullet}}

% ── Кодын блок ──────────────────────────────────────────────────────────────
\usepackage{fvextra}
\usepackage{framed}
\fvset{
  breaklines=true,
  breakanywhere=true,
  breaksymbolleft={},
  breakindent=0pt,
  fontsize=\footnotesize,
  commandchars=\\\{\},
}
% pandoc-ийн highlighting макронууд (highlighting.tex файлаас)
\usepackage{color}
\usepackage{fancyvrb}
\newcommand{\VerbBar}{|}
\newcommand{\VERB}{\Verb[commandchars=\\\{\}]}
\DefineVerbatimEnvironment{Highlighting}{Verbatim}{commandchars=\\\{\}}
% Add ',fontsize=\small' for more characters per line
\usepackage{framed}
\definecolor{shadecolor}{RGB}{248,248,248}
\newenvironment{Shaded}{\begin{snugshade}}{\end{snugshade}}
\newcommand{\AlertTok}[1]{\textcolor[rgb]{0.94,0.16,0.16}{#1}}
\newcommand{\AnnotationTok}[1]{\textcolor[rgb]{0.56,0.35,0.01}{\textbf{\textit{#1}}}}
\newcommand{\AttributeTok}[1]{\textcolor[rgb]{0.13,0.29,0.53}{#1}}
\newcommand{\BaseNTok}[1]{\textcolor[rgb]{0.00,0.00,0.81}{#1}}
\newcommand{\BuiltInTok}[1]{#1}
\newcommand{\CharTok}[1]{\textcolor[rgb]{0.31,0.60,0.02}{#1}}
\newcommand{\CommentTok}[1]{\textcolor[rgb]{0.56,0.35,0.01}{\textit{#1}}}
\newcommand{\CommentVarTok}[1]{\textcolor[rgb]{0.56,0.35,0.01}{\textbf{\textit{#1}}}}
\newcommand{\ConstantTok}[1]{\textcolor[rgb]{0.56,0.35,0.01}{#1}}
\newcommand{\ControlFlowTok}[1]{\textcolor[rgb]{0.13,0.29,0.53}{\textbf{#1}}}
\newcommand{\DataTypeTok}[1]{\textcolor[rgb]{0.13,0.29,0.53}{#1}}
\newcommand{\DecValTok}[1]{\textcolor[rgb]{0.00,0.00,0.81}{#1}}
\newcommand{\DocumentationTok}[1]{\textcolor[rgb]{0.56,0.35,0.01}{\textbf{\textit{#1}}}}
\newcommand{\ErrorTok}[1]{\textcolor[rgb]{0.64,0.00,0.00}{\textbf{#1}}}
\newcommand{\ExtensionTok}[1]{#1}
\newcommand{\FloatTok}[1]{\textcolor[rgb]{0.00,0.00,0.81}{#1}}
\newcommand{\FunctionTok}[1]{\textcolor[rgb]{0.13,0.29,0.53}{\textbf{#1}}}
\newcommand{\ImportTok}[1]{#1}
\newcommand{\InformationTok}[1]{\textcolor[rgb]{0.56,0.35,0.01}{\textbf{\textit{#1}}}}
\newcommand{\KeywordTok}[1]{\textcolor[rgb]{0.13,0.29,0.53}{\textbf{#1}}}
\newcommand{\NormalTok}[1]{#1}
\newcommand{\OperatorTok}[1]{\textcolor[rgb]{0.81,0.36,0.00}{\textbf{#1}}}
\newcommand{\OtherTok}[1]{\textcolor[rgb]{0.56,0.35,0.01}{#1}}
\newcommand{\PreprocessorTok}[1]{\textcolor[rgb]{0.56,0.35,0.01}{\textit{#1}}}
\newcommand{\RegionMarkerTok}[1]{#1}
\newcommand{\SpecialCharTok}[1]{\textcolor[rgb]{0.81,0.36,0.00}{\textbf{#1}}}
\newcommand{\SpecialStringTok}[1]{\textcolor[rgb]{0.31,0.60,0.02}{#1}}
\newcommand{\StringTok}[1]{\textcolor[rgb]{0.31,0.60,0.02}{#1}}
\newcommand{\VariableTok}[1]{\textcolor[rgb]{0.00,0.00,0.00}{#1}}
\newcommand{\VerbatimStringTok}[1]{\textcolor[rgb]{0.31,0.60,0.02}{#1}}
\newcommand{\WarningTok}[1]{\textcolor[rgb]{0.56,0.35,0.01}{\textbf{\textit{#1}}}}
\setlength{\emergencystretch}{3em} % prevent overfull lines
\providecommand{\tightlist}{%
  \setlength{\itemsep}{0pt}\setlength{\parskip}{0pt}}
\setcounter{secnumdepth}{-\maxdimen} % remove section numbering
\ifLuaTeX
  \usepackage{selnolig}  % disable illegal ligatures
\fi
\IfFileExists{bookmark.sty}{\usepackage{bookmark}}{\usepackage{hyperref}}
\IfFileExists{xurl.sty}{\usepackage{xurl}}{} % add URL line breaks if available
\urlstyle{same}
\hypersetup{
  hidelinks,
  pdfcreator={LaTeX via pandoc}}

\author{}
\date{}


% Кодын хүрээг цэвэрлэх
\definecolor{shadecolor}{HTML}{F7F8FA}
\setlength{\FrameSep}{5pt}
\setlength{\OuterFrameSep}{4pt}


% ── Тэмдэглэл, анхааруулгын хайрцаг ─────────────────────────────────────────
\usepackage[most]{tcolorbox}
\tcbuselibrary{breakable}

\newtcolorbox{cncnote}{
  breakable, enhanced,
  colback=cnclight!30, colframe=cncmid,
  boxrule=0pt, leftrule=2.5pt,
  arc=0pt, outer arc=0pt,
  left=7pt, right=7pt, top=5pt, bottom=5pt,
  before skip=7pt, after skip=7pt,
}

% Markdown-ийн `>` блокийг хайрцаг болгоно
\renewenvironment{quote}{\begin{cncnote}\small}{\end{cncnote}}

% ── LaTeX-ийн бэлэн нэрсийг монголчлох ─────────────────────────────────────
\renewcommand{\contentsname}{Агуулга}
\renewcommand{\chaptername}{Бүлэг}
\renewcommand{\appendixname}{Хавсралт}
\renewcommand{\tablename}{Хүснэгт}
\renewcommand{\figurename}{Зураг}
\renewcommand{\listtablename}{Хүснэгтийн жагсаалт}
\renewcommand{\listfigurename}{Зургийн жагсаалт}
\renewcommand{\bibname}{Ном зүй}
\renewcommand{\indexname}{Хэлхээ бүртгэл}
\renewcommand{\partname}{Хэсэг}

% ── Гарчиг ──────────────────────────────────────────────────────────────────
\usepackage{titlesec}

\titleformat{\chapter}[display]
  {\normalfont\sffamily\bfseries\color{cncblue}}
  {\raggedright\normalsize\sffamily\color{cncmid}\MakeUppercase{\chaptertitlename}\ \thechapter}
  {6pt}
  {\Huge\raggedright}
  [\vspace{2pt}{\color{cncrule}\titlerule[1.2pt]}]
\titlespacing*{\chapter}{0pt}{-18pt}{22pt}

\titleformat{\section}
  {\normalfont\sffamily\Large\bfseries\color{cncblue}}{\thesection}{0.7em}{}
\titleformat{\subsection}
  {\normalfont\sffamily\large\bfseries\color{cncmid}}{\thesubsection}{0.6em}{}
\titleformat{\subsubsection}
  {\normalfont\sffamily\normalsize\bfseries\color{cncmid}}{\thesubsubsection}{0.5em}{}
\titlespacing*{\section}{0pt}{14pt}{5pt}
\titlespacing*{\subsection}{0pt}{11pt}{4pt}
\titlespacing*{\subsubsection}{0pt}{9pt}{3pt}

% Заавар өөрөө '1. Зорилго', 'Алхам 6' гэж дугаарласан тул LaTeX-ийн
% давхар дугаарлалтыг унтраана (бүлэг л дугаарлагдана).
\setcounter{secnumdepth}{0}
\setcounter{tocdepth}{1}

% ── Толгой, хөл ─────────────────────────────────────────────────────────────
\usepackage{fancyhdr}
\pagestyle{fancy}
% DejaVu Sans-ийн мөрийн өндөр geometry-гийн тооцооноос өндөр тул fancyhdr
% анхааруулга өгдөг. Энд ил зааж, дээд захыг тэр хэмжээгээр нөхнө.
\setlength{\headheight}{22pt}
\addtolength{\topmargin}{-4pt}
\fancyhf{}
\renewcommand{\headrulewidth}{0.4pt}
\renewcommand{\headrule}{\hbox to\headwidth{\color{cncrule}\leaders\hrule height \headrulewidth\hfill}}
% Бүлгийн гарчиг урт (жишээ: "Лаб 1 — Хоёр давхаргат платформ ба ирмэг–үүлний
% гүүр") тул зүүн талыг зөвхөн хичээлийн код болгож давхцахаас сэргийлнэ.
\fancyhead[L]{\footnotesize\sffamily\color{cncmid}CNC302}
\fancyhead[R]{\footnotesize\sffamily\color{cncmid}\nouppercase{\leftmark}}
\fancyfoot[C]{\small\thepage}
\fancypagestyle{plain}{\fancyhf{}\fancyfoot[C]{\small\thepage}\renewcommand{\headrulewidth}{0pt}}
\renewcommand{\chaptermark}[1]{\markboth{#1}{}}

% ── Холбоос ─────────────────────────────────────────────────────────────────
\usepackage{hyperref}
\hypersetup{
  unicode=true,
  bookmarks=true,
  bookmarksnumbered=true,
  colorlinks=true,
  linkcolor=cncmid,
  urlcolor=cncmid,
  citecolor=cncmid,
  pdftitle={CNC302 — Лабораторийн ажлын гарын авлага},
  pdfauthor={С.Батдалай, ШУТИС МХТС Холбооны тэнхим},
  pdfsubject={Cloud-Native IoT Platforms and Applications},
  pdfkeywords={IoT, MQTT, Kubernetes, Edge AI, Raspberry Pi},
}
\urlstyle{same}

% Урт URL-ийг таслах
\usepackage{xurl}

% ── Тусламжийн макронууд ────────────────────────────────────────────────────
\newcommand{\labmeta}[5]{%
  \begin{center}
  \begin{tabular}{@{}>{\bfseries\sffamily\color{cncmid}}l l@{}}
    \toprule
    7 хоног & #1 \\
    Хугацаа & #2 \\
    Суралцахуйн үр дүн & #3 \\
    Үнэлгээ & #4 \\
    Гол хэмжилт & #5 \\
    \bottomrule
  \end{tabular}
  \end{center}
  \vspace{4pt}
}

% Хүснэгтийг жижигрүүлэх орчин (өргөн хүснэгтэд)
\newenvironment{widetable}{\begingroup\scriptsize}{\endgroup}

\newcommand{\keyterm}[1]{\textbf{\textcolor{cncblue}{#1}}}
\newcommand{\file}[1]{\texttt{\small #1}}

% Markdown-ийн шугам (---) -ийг зөөлөн зураас болгоно
\providecommand{\mdrule}{\par\vspace{3pt}{\color{cncrule}\hrule height 0.4pt}\par\vspace{6pt}}

% pandoc-ийн шаарддаг макронууд
\providecommand{\tightlist}{\setlength{\itemsep}{0pt}\setlength{\parskip}{0pt}}
\providecommand{\pandocbounded}[1]{#1}
\providecommand{\passthrough}[1]{#1}

% ── Кодын орчны эцсийн тодорхойлолт ─────────────────────────────────────────
% ЧУХАЛ: энэ хэсэг preamble-ийн ХАМГИЙН СҮҮЛД байх ёстой. tcolorbox зэрэг
% багцууд verbatim орчныг дахин тодорхойлдог тул эрт тавибал үйлчлэхгүй.
\AtBeginDocument{%
  \DefineVerbatimEnvironment{Highlighting}{Verbatim}{%
    commandchars=\\\{\},breaklines=true,breakanywhere=true,%
    breaksymbolleft={},fontsize=\scriptsize}%
  \DefineVerbatimEnvironment{verbatim}{Verbatim}{%
    breaklines=true,breakanywhere=true,breaksymbolleft={},fontsize=\scriptsize}%
}
